\VerdierPhysicalPageAnchor{0022}%
à homotopie près par le complexe $Y^{\bullet}$. Les objets de cohomologie
$\mathrm R^nF(Y^{\bullet})$ du complexe
$\int F(I^{\bullet\bullet})$ sont donc définis à isomorphisme canonique près
et on constate qu'ils dépendent fonctoriellement du complexe $Y^{\bullet}$.
Les foncteurs $\mathrm R^nF$ obtenus sont les
\emph{hyperdérivés à droite} de $F$.

\VerdierSemanticAnchor{intro-2-spectral-sequences}
On obtient de plus deux suites spectrales fonctorielles en $Y^{\bullet}$ dont
les termes initiaux sont :

\VerdierSemanticAnchor{intro-2-spectral-sequence-one}
\begin{equation}
\mathrm I_2^{p,q}=\mathrm H^p\bigl(\mathrm R^qF(Y^{\bullet})\bigr)
\tag{*}
\end{equation}

\VerdierSemanticAnchor{intro-2-spectral-sequence-two}
\begin{equation}
\mathrm{II}_2^{p,q}=\mathrm R^pF\bigl(\mathrm H^q(Y^{\bullet})\bigr),
\tag{**}
\end{equation}
et dont les termes infinis sont les gradués associés aux objets
$\mathrm R^nF(Y^{\bullet})$ convenablement filtrés. Bornons-nous, pour
simplifier, au cas où les composants des complexes envisagés sont nuls en
degré suffisamment petit. On constate alors que les suites spectrales (*) et
(**) convergent. Un morphisme $m:Y^{\bullet}\to Z^{\bullet}$ de complexes
induit des morphismes :

\VerdierSemanticAnchor{intro-2-hyperderived-morphism}
\[
\mathrm R^nF(m):\mathrm R^nF(Y^{\bullet})
\longrightarrow \mathrm R^nF(Z^{\bullet})
\]
et des morphismes des suites spectrales correspondantes.

\VerdierSemanticAnchor{intro-2-quasi-isomorphism-invariance}
En particulier, \emph{lorsque $m$ est un quasi-isomorphisme}, la suite
spectrale (**) montre que \emph{les morphismes $\mathrm R^nF(m)$ sont des
isomorphismes}. Ainsi les objets $\mathrm R^nF(Y^{\bullet})$ sont déterminés
dès que l'on connaît le complexe $Y^{\bullet}$ à ``quasi-isomorphisme près''.
Comme, par ailleurs, on a indiqué que les procédés usuels de l'algèbre
homologique ne fournissent le plus souvent que des complexes déterminés à
quasi-isomorphisme près, il apparaît que les objets naturels qu'étudie
l'hyperhomologie sont les ``classes de complexes à quasi-isomorphisme près''.

\VerdierSemanticAnchor{intro-3}
\noindent\textbf{3.}\quad
La démarche que nous suivons s'impose alors naturellement. On introduit, pour
toute catégorie abélienne $\mathcal A$, la catégorie
$\operatorname{comp}(\mathcal A)$ des complexes d'objets de $\mathcal A$.
Puis on construit à partir de $\operatorname{comp}(\mathcal A)$ une nouvelle
catégorie en rendant formellement inversibles les quasi-isomorphismes de
$\operatorname{comp}(\mathcal A)$.

\VerdierSemanticAnchor{intro-3-derived-category}
On obtient ainsi une catégorie $\mathbf D(\mathcal A)$ :
\emph{la ``catégorie dérivée'' de la catégorie abélienne $\mathcal A$}.
Par des procédés de résolution à droite qu'il importe peu de préciser dans
cette introduction, on associe à un foncteur additif entre deux catégories
abéliennes $F:\mathcal A\to\mathcal B$, un foncteur
$\mathrm RF:\mathbf D(\mathcal A)\to\mathbf D(\mathcal B)$ (qui n'est défini
le plus souvent que sur une sous-catégorie convenable
$\mathbf D^*(\mathcal A)\subset\mathbf D(\mathcal A)$ obtenue en imposant des
limitations aux degrés des complexes envisagés) appelé le
\emph{foncteur dérivé total à droite} de $F$.

\VerdierSemanticAnchor{intro-3-total-right-derived-functor}
Pour tout objet $X$ de $\mathbf D^*(\mathcal A)$, l'objet $\mathrm RF(X)$ de
$\mathbf D(\mathcal B)$ est donc un complexe d'objets de $\mathcal B$ défini à
``quasi-isomorphisme près'', et par suite, les objets de cohomologie
$\mathrm R^nF(X)$ de ce complexe sont bien définis. Les foncteurs
$X\longmapsto\mathrm R^nF(X)$ sont les ``hyperdérivés à droite'' du foncteur
$F$, définis ici sans supposer nécessairement que
